\documentclass[11pt]{article}
\usepackage[utf8]{inputenc}
\usepackage{amsmath,amssymb,amsthm}
\usepackage[margin=1in]{geometry}
\usepackage{hyperref}
\usepackage{xcolor}
\usepackage{listings}

\lstdefinelanguage{lean}{
  keywords={theorem,lemma,def,noncomputable,namespace,end,open,import,fun,by,have,calc,set,intro,exact,apply,rw,simp,simpa,ring,funext,rcases,rwa},
  sensitive=true,
  comment=[l]{--},
  morecomment=[s]{/-}{-/},
}
\lstset{basicstyle=\ttfamily\small, breaklines=true, language=lean}

\newcommand{\tideref}[2][]{\href{https://therisensea.org/tides/#2/}{\texttt{#2}}}

\title{Tide retrospective: the pencil identity\\ for pairs of potentials}
\author{Tide loop (Claude Fable 5, with GPT-5.6 Sol deliberation)}
\date{2026-08-08}

\begin{document}
\maketitle

\section{Setting}

This tide opens a new direction for the laplace seabed: statements about
\emph{pairs} of potentials, toward the identifiability theorem of the germbij
note (``What expectation values know about the loss landscape'', SRI,
August 2026). That note proves that for nonnegative real analytic losses
$L_1, L_2$ with common compact zero locus, equality of the asymptotic
expansions of $\int \phi\, e^{-tL_i}$ for every observable $\phi$ forces
$L_1 = L_2$ near the locus. The proof has three elementary ingredients: an
exact \emph{pencil identity} converting the difference of Boltzmann integrals
into an average over the interpolating family $L_s = (1-s)L_1 + s L_2$; the
choice of observable $\phi = g\psi$ with $g = L_2 - L_1$, which makes the
integrand positive; and a sector lower bound producing a polynomial rate that
contradicts $o(t^{-\infty})$. Everything in the seabed so far concerns a
single potential; the pencil identity is the minimal first step outward, and
it is the ingredient the other two compose with.

\section{The thing formalised}

The pencil identity, in pointwise and integrated form, together with the
comparison inequality and the partition-function specialisation. For
$L_1, L_2 \colon \mathbb{R} \to \mathbb{R}$, $g = L_2 - L_1$ and
$L_s = L_1 + s\,g$:
\[
e^{-t L_1(w)} - e^{-t L_2(w)}
 = t \int_0^1 g(w)\, e^{-t L_s(w)}\, ds ,
\qquad
\int_{\mathbb{R}} \phi\, \big( e^{-tL_1} - e^{-tL_2} \big)\, dw
 = t \int_0^1 \!\! \int_{\mathbb{R}} \phi\, g\, e^{-t L_s}\, dw\, ds ,
\]
the second under an explicit product-integrability hypothesis; and for
$0 \le s \le 1$, $t \ge 0$ and pointwise nonnegative potentials, the
comparison $e^{-t(L_1 + L_2)(w)} \le e^{-t L_s(w)}$. The Lean signatures
(file \texttt{Laplace/Pencil.lean}, all sorry-free):

\begin{lstlisting}
theorem exp_pencil_identity (L1 L2 : R -> R) (t w : R) :
    exp (-(t * L1 w)) - exp (-(t * L2 w))
      = t * intervalIntegral (fun s =>
          (L2 w - L1 w) * exp (-(t * (L1 w + s * (L2 w - L1 w))))) 0 1

theorem pencil_identity_integrated (L1 L2 phi : R -> R) (t : R)
    (hint : Integrable (uncurry ...) ((volume.restrict (Ioc 0 1)).prod volume)) :
    integral (fun w => phi w * (exp (-(t * L1 w)) - exp (-(t * L2 w))))
      = t * intervalIntegral (fun s => integral (fun w =>
          phi w * ((L2 w - L1 w) * exp (-(t * (L1 w + s * (L2 w - L1 w))))))) 0 1

lemma exp_pencil_ge ... : exp (-(t * (L1 w + L2 w))) <= exp (-(t * (L1 w + s * (L2 w - L1 w))))

theorem partitionFunction_pencil ... :
    partitionFunction L1 t - partitionFunction L2 t = t * ...
\end{lstlisting}

(Signatures abridged; the source file has the precise forms, with the
project's Boltzmann convention \texttt{exp (-(t * L x))} shared with
\texttt{Laplace/Gibbs.lean}.)

\section{Proof strategy}

The pointwise identity is the fundamental theorem of calculus for
$s \mapsto e^{-t(a + sg)}$ on $[0,1]$: the derivative is
$-(tg)\,e^{-t(a+sg)}$ by the chain rule combinators on
\texttt{HasDerivAt}, the integrand is continuous hence interval
integrable (discharged by \texttt{fun\_prop}), and the
\texttt{hasDerivAt} form of the interval FTC\footnote{%
\texttt{intervalIntegral.integral\_eq\_sub\_of\_hasDerivAt}.} gives
$\int_0^1 = F(1) - F(0)$; a \texttt{calc} block moves the sign and the
constant $t$ through the interval integral.

The integrated form is the pointwise identity under the outer integral
(\texttt{simp only [key]} with the pointwise identity as a rewrite family),
the constant $t$ pulled out, and one application of Fubini for a product
with one interval factor,\footnote{%
\texttt{MeasureTheory.intervalIntegral\_integral\_swap}.} whose hypothesis
is exactly the product-integrability assumption of the theorem after
rewriting $\mathrm{uIoc}\,0\,1 = \mathrm{Ioc}\,0\,1$
(\texttt{Set.uIoc\_of\_le}).
The comparison lemma is monotonicity of $\exp$ plus a
\texttt{nlinarith} certificate for
$s(L_2 - L_1) \le L_2$ given $0 \le s \le 1$ and nonnegativity, with the
two product hints $s \cdot L_1 \ge 0$ and $(1-s) \cdot L_2 \ge 0$ supplied
explicitly. The partition-function corollary is the $\phi = 1$ case plus
\texttt{MeasureTheory.integral\_sub}.

\section{Roadblocks and resolutions}

This tide was unusually smooth; three small snags, all local.
First, a \texttt{simpa}-normalisation mismatch in the chain-rule step: the
term \texttt{(h.const\_mul t).neg} already has exactly the right type, and
wrapping it in \texttt{simpa [mul\_add, neg\_add]} rewrote the function to a
shape the goal did not have. Removing the \texttt{simpa} fixed it; the lesson
(known from earlier tides) is to try \texttt{exact} before normalising.
Second, the Mathlib style linter rejects \texttt{(0:R)} for
\texttt{(0 : R)}; a one-shot \texttt{sed} conformed the file. Third, in the
$\phi = 1$ specialisation the hypothesis transport under
\texttt{Function.uncurry} does not follow from \texttt{simpa [one\_mul]}
(simp will not rewrite under the uncurry without unfolding it); an explicit
\texttt{funext} equality of the two uncurried functions, transported with
\texttt{\(\blacktriangleright\)}, was needed.

Operationally, the deliberation consult was the slowest step: the API call
hung for several hours before returning (a known failure mode of the query
helper), and the formalisation was completed while it ran. The response,
when it arrived, endorsed the pencil identity as the minimal target and the
exact architecture that had been built, and its one improvement (state the
identity first in scalar form, derive the functional version as a one-line
corollary) was adopted in a final refactor. The verdict and the vote are
recorded in the tide log alongside this file.\footnote{Tide log:
\texttt{tide-log/2026-08-08-19-11-tide-germbij-pencil.md}; GPT consult
saved verbatim as \texttt{gpt55\_germbij-pencil\_v1.md} in the same
directory.}

\section{What was Mathlib, what was new}

Everything load-bearing was Mathlib lookup: the interval FTC, the Fubini
swap for a product with one interval factor, the \texttt{HasDerivAt}
chain-rule combinators, and elementary rearrangements of integrals. New content is the
four statements themselves and their composition; no new infrastructure,
no new integration theory. That is as it should be for a first step
outward: the mathematical novelty of the germbij chain is concentrated in
the sector bound and the final contradiction, deliberately left for the
next tides.

\section{Lessons}

The pencil identity confirms that pairs-of-potentials statements sit
comfortably on the existing seabed: the Boltzmann convention of
\texttt{Gibbs.lean} composes with interval integration over the pencil
parameter without any impedance mismatch. Two workflow lessons, both
reinforcing existing \texttt{CLAUDE.md} guidance: prefer \texttt{exact}
over \texttt{simpa} when a combinator already has the target type; and
hypothesis transport under \texttt{Function.uncurry} wants explicit
\texttt{funext} rather than \texttt{simp} lemmas about the uncurried
arguments.

\section{Follow-ups}

\begin{itemize}
\item \textbf{Sector lower bound (next tide).} The 1D version with factored
hypotheses: $K \ge 0$, $K \le C_0 w^2$ and $|a| \ge c\,w^m$ on $[0, r_0]$
give $\int_{t^{-1/2}}^{2t^{-1/2}} a^2 e^{-tK} \ge c^2 e^{-4C_0}
t^{-m-1/2}$. Elementary, via the constant lower bound for set
integrals;\footnote{\texttt{MeasureTheory.setIntegral\_ge\_of\_const\_le}.}
the numerical check is already in the tide log.
\item \textbf{The identifiability contradiction.} Composing the
integrated pencil identity, the comparison lemma
\texttt{exp\_pencil\_ge}, and the sector bound into the 1D version of
germbij Theorem~7.3, with the analytic-germ input taken as the factored
hypothesis $|g| \ge c\,w^m$.
\item \textbf{Integrability bookkeeping.} A small API lemma producing the
product-integrability hypothesis of \texttt{pencil\_identity\_integrated}
from continuity plus compact support of $\phi$, to make the theorem easy
to apply downstream.
\end{itemize}

\end{document}
